\section{\label{sec:experiments}Experiments}

\subsection{Training configuration}

We train on 2 NVIDIA H200 GPUs using distributed data parallel (DDP) with bfloat16 mixed-precision (model forward pass
in bfloat16, all loss computation in float32 for numerical stability). The complex-valued U-Net uses a base channel
count of 40, channel multipliers $(1, 2, 4, 8)$, 3 residual blocks per level, and 8-head complex self-attention at the
$8\times$-downsampled resolution. We use AdamW~\cite{Loshchilov2019AdamW} with learning rate $10^{-4}$, weight decay
$5 \times 10^{-5}$, and a schedule combining 10-epoch linear warmup with cosine decay to a minimum of $5 \times
10^{-6}$. The global batch size is 160 across both GPUs. An exponential moving average (EMA) of model weights with
decay 0.999 is maintained and used for all evaluation.

The training curriculum is: Phase~A for 30 epochs (FM only), Phase~B from epoch 30 to 100 (add Helmholtz residual
$\lambda_R = 1.0$ with 20-epoch linear warmup), and Phase~C from epoch 100 onward (add phase loss $\lambda_\phi =
0.1$ and S-parameter loss $\lambda_S = 0.1$, each with 30-epoch warmups). S-parameters are extracted using the
projection method. Physics losses are time-gated to samples with $t \geq 0.5$, and the Helmholtz residual is
additionally weighted by $t^4$ to upweight clean, high-$t$ reconstructions. Interface masking with threshold $\tau =
1.0$ excludes dielectric boundary pixels from the residual. Conflict-free gradient surgery
(ConFIG)~\cite{ConFIG2025} activates at epoch 100 when all four losses are simultaneously active.
Time is sampled from a 50/50 mixture of uniform and logit-normal distributions. Inference uses Heun integration with
30 steps on a linear time grid. Training runs for 400 epochs total.

\subsection{Evaluation metrics}

\paragraph{Field quality.}
We report peak signal-to-noise ratio (PSNR) on the field magnitude $|E_z|$ and relative $L_2$ amplitude error, both
computed on the held-out test set after global phase alignment. Phase error is reported as the root-mean-square angular
deviation (in degrees) between predicted and ground-truth unit phasors, weighted by $|E_z^\text{true}|^2$ and gated by
an amplitude threshold to suppress noise in low-signal regions.

\paragraph{Physics compliance.}
The Helmholtz residual ratio (Sec.~\ref{sec:method}) with interface masking, computed on predicted fields and
normalized by the driving-term scale. We report this relative to the ground-truth floor: a ratio of $1.0\times$~GT
indicates the model satisfies the Helmholtz equation as well as the FDTD reference solution itself.

\paragraph{S-parameter accuracy.}
Mean absolute error in $|S_{ji}|$ and phase error (degrees) for each output port, broken down by device type. The
input port is excluded from evaluation since the total-field projection at the excited port conflates incident and
scattered waves.

\paragraph{Computational cost.}
Wall-clock inference time per sample (single and batched) on a single GPU, compared to FDTD simulation time. We also
report model parameter count and total training GPU-hours.

\subsection{Baselines and ablations}

We evaluate the following configurations to isolate the contribution of each design choice:

\begin{enumerate}
    \item \textbf{Complex U-Net + full curriculum} (ours): The complete model as described in Sec.~\ref{sec:method}.
    \item \textbf{Real-valued U-Net}: Same architecture and training procedure, but with standard real convolutions
          applied independently to $E_r$ and $E_i$, standard ReLU, and real group normalization. Tests whether
          complex-valued inductive bias improves phase accuracy.
    \item \textbf{FM only}: Flow matching loss without Helmholtz residual, phase, or S-parameter losses. Tests
          whether physics-constrained training improves field quality beyond what pure data-driven learning achieves.
    \item \textbf{No curriculum}: All losses active from epoch 0 with their full target weights (no phased
          introduction, no warmup ramps). Tests whether the multi-loss curriculum is necessary for stable training.
    \item \textbf{No interface masking}: Helmholtz residual computed at all interior pixels including dielectric
          boundaries. Tests whether interface masking improves training dynamics and the usefulness of the residual
          as a compliance metric.
\end{enumerate}

\subsection{Results}

% -----------------------------------------------------------------------
% PLACEHOLDER: Fill in after training runs complete.
% -----------------------------------------------------------------------

\subsubsection{Field prediction quality}

% Table: PSNR, amplitude error, phase error per device type
% Compare complex vs real UNet, FM-only vs full curriculum

\begin{table*}[t]
\centering
\caption{Field prediction quality on the held-out test set (EMA model, 30-step Heun integration). PSNR is computed on
$|E_z|$; phase error is amplitude-weighted RMS in degrees. Best result per column in bold.}
\label{tab:field_quality}
\begin{tabular}{lccccc}
\hline\hline
 & \multicolumn{5}{c}{PSNR (dB) $\uparrow$} \\
\cline{2-6}
Method & S-bend & Y-branch & MMI & Dir.\ coupler & Mean \\
\hline
FM only           & -- & -- & -- & -- & -- \\
Real U-Net + curriculum & -- & -- & -- & -- & -- \\
Complex U-Net + curriculum (ours) & -- & -- & -- & -- & -- \\
\hline\hline
\end{tabular}
\end{table*}

\subsubsection{Physics compliance}

% Table or figure: Helmholtz residual ratio (with and without interface masking)
% Show that interface masking drops GT floor from ~0.17 to ~0.01-0.03
% Show model's residual ratio relative to GT floor

\subsubsection{S-parameter accuracy}

% Table: |S| magnitude error and phase error per device type and port
% Compare project vs modal extraction
% Show that S-param loss improves accuracy vs FM-only

\begin{table}[t]
\centering
\caption{S-parameter accuracy on the test set. Magnitude error is mean $||S^\text{pred}| - |S^\text{true}||$; phase
error is RMS angular difference in degrees, evaluated on output ports with $|S^\text{true}| > 0.01$.}
\label{tab:sparam}
\begin{tabular}{lcc}
\hline\hline
Device & $\Delta|S|$ $\downarrow$ & $\Delta\angle S$ (deg) $\downarrow$ \\
\hline
S-bend          & -- & -- \\
Y-branch        & -- & -- \\
MMI             & -- & -- \\
Dir.\ coupler   & -- & -- \\
\hline
Mean            & -- & -- \\
\hline\hline
\end{tabular}
\end{table}

\subsubsection{Ablation study}

% Table: All ablation variants with PSNR, residual ratio, S-param error
% Key findings:
% - Complex vs real: complex improves phase accuracy by X
% - FM-only vs full: physics losses improve residual ratio by Y
% - No curriculum vs curriculum: curriculum stabilizes training, improves final metrics
% - No interface masking: residual floor is artificially high, noisy gradients

\begin{table}[t]
\centering
\caption{Ablation study (test-set averages across all device types).}
\label{tab:ablation}
\begin{tabular}{lccc}
\hline\hline
Variant & PSNR $\uparrow$ & Res.\ ratio $\downarrow$ & $\Delta|S|$ $\downarrow$ \\
\hline
Ours (full)           & -- & -- & -- \\
Real-valued U-Net     & -- & -- & -- \\
FM only               & -- & -- & -- \\
No curriculum         & -- & -- & -- \\
No interface mask     & -- & -- & -- \\
\hline\hline
\end{tabular}
\end{table}

\subsubsection{Computational cost}

% Table: inference time vs FDTD
% Single sample, batched (e.g. batch=32), vs FDTD per sim
% Parameter count
% Training GPU-hours

\begin{table}[t]
\centering
\caption{Computational cost comparison.}
\label{tab:cost}
\begin{tabular}{lc}
\hline\hline
 & Time \\
\hline
FDTD (Meep, single sim)           & $\sim$\,2--5\,min \\
Ours (single sample, 1 GPU)       & -- \\
Ours (batch of 32, 1 GPU)         & -- \\
\hline
Model parameters                   & -- \\
Training (2$\times$ H200, 400 ep) & --\,GPU-hours \\
\hline\hline
\end{tabular}
\end{table}

\subsubsection{Qualitative examples}

% Figure: Side-by-side GT vs predicted for representative devices
% Rows: amplitude |Ez|, phase angle(Ez), pointwise error
% Columns: one device per column (S-bend, Y-branch, MMI, dir coupler)
% Caption: note that phase is accurately captured in high-amplitude regions

% Figure: Helmholtz residual maps with and without interface masking
% Show that interface pixels dominate residual without masking

% Figure: Training loss curves showing Phase A/B/C transitions
% Show FM loss, residual loss, phase loss, sparam loss vs epoch
